\section{Technical Overview}
\label{sec:overview}

\methodname consists of three components.  First, it reorganizes a collection
of heterogeneous matrices into one dense vector without materializing
per-matrix padding.  Second, it converts independent MLE claims on the
matrices into one inner-product claim on this vector, using only committed
shape information.  Finally, it proves the resulting inner product directly
from the native low-bit-width representation.

\subsection{Splitting and Recombining the Matrices}

Consider a collection of matrices
\[
    \{\bm C_\tau\}_{\tau=1}^{t}
\]
containing \(N\) entries in total, and let \(K\) be the next power of two
greater than or equal to \(N\).
We will show that each matrix can be partitioned into at most
\(\log_2 K^2\) aligned submatrices whose row and column dimensions are
powers of two.  Across the collection, this produces:
\[
    T\le t\log_2 K^2
\]
submatrices.

We flatten these submatrices in column-major order, sort them by non-increasing
size, and concatenate them without gaps.  The resulting dense vector
\[
    \bm a\in\Field^N
\]
contains every matrix entry exactly once.  We call this representation the
\emph{bucketed vector assembly} (BVA) of the matrix collection.  Its logical
PCS domain has length \(K\), but the final \(K-N\) zero entries remain
implicit.  A commitment to \(\bm a\), together with a consistency proof for
the splitting, therefore binds the entire matrix collection through a single
data commitment.

\subsection{MLE-Claim Conversion with Hidden Shape Encodings}

Suppose that we wish to prove the independent MLE claims
\[
    u_\tau
    =
    \Eval(
        \bm C_\tau,
        \bm x_{\tau,l},
        \bm x_{\tau,r}
    ),
    \qquad \tau\in[t],
\]
where \(\Eval(\bm C_\tau,\bm x_{\tau,l},\bm x_{\tau,r})\) denotes the
matrix MLE evaluation at the left and right points
\((\bm x_{\tau,l},\bm x_{\tau,r})\).

After sampling a batching challenge \(\omega\), we aggregate the claims as
\[
    c
    :=
    \sum_{\tau=1}^{t}
        \omega^{\tau-1}u_\tau.
\]
Each power-of-two submatrix corresponds to one segment of \(\bm a\).
Because the submatrix is aligned, restricting the tensor evaluation vector of
its source matrix to this segment yields
\[
    \eta_j\tensorspan{\bm z_j},
\]
where \(\bm z_j\) contains the coordinates that vary within the submatrix and
\(\eta_j\) contains the batching coefficient and the coordinates fixed by its
offset.  Concatenating these tensor-structured segments in the same order as
the BVA produces a semi-structured vector \(\bm b\) satisfying
\[
    c
    =
    \sum_{j=1}^{T}
        \eta_j
        \tensorproj{
            \vectorize{\bm C^{(j)}}
        }{\bm z_j}
    =
    \langle\bm a,\bm b\rangle.
\]

The conversion from
\(\{(\bm x_{\tau,l},\bm x_{\tau,r})\}_{\tau=1}^{t}\)
to \(\{(\eta_j,\bm z_j)\}_{j=1}^{T}\) depends only on the source-matrix
assignment, the aligned offsets, and the dimensions of the submatrices.
We encode this shape information injectively and commit to the encoding.
Lookup arguments derive the local points \(\bm z_j\) from the original
evaluation points, while row-wise grand-product arguments derive the scales
\(\eta_j\).  Zero-knowledge consistency proofs bind this conversion to the
committed encoding without revealing the matrix shapes.

\subsection{Packed Low-Bit-Width Inner-Product Proof}

It remains to prove
\[
    c=\langle\bm a,\bm b\rangle,
\]
where \(\bm a\) is stored using \(b\)-bit values and \(\bm b\) is a
semi-structured field vector.  Choose
\[
    \kappa=\lfloor |\Field|/b \rfloor
\]
and partition the logical length-\(K\) vectors into \(\kappa\) equal-length
segments
\(\{\bm a_{(s)}\}_{s=1}^{\kappa}\) and
\(\{\bm b_{(s)}\}_{s=1}^{\kappa}\).  Define the Laurent-polynomial
encodings
\[
    \bm g_a(Y)
    :=
    \sum_{s=1}^{\kappa}
        Y^{s-1}\bm a_{(s)},
    \qquad
    \bm g_b(Y)
    :=
    \sum_{s=1}^{\kappa}
        Y^{-(s-1)}\bm b_{(s)}.
\]
Their inner product has the form
\begin{equation}
\label{eq:packed-ip-overview}
\left\langle
    \bm g_a(Y),\bm g_b(Y)
\right\rangle
=
c+
\sum_{s=1}^{\kappa-1}
\left(
    c_{+,s}Y^s
    +
    c_{-,s}Y^{-s}
\right),
\end{equation}
because the constant coefficient is
\[
    \sum_{s=1}^{\kappa}
        \langle\bm a_{(s)},\bm b_{(s)}\rangle
    =
    \langle\bm a,\bm b\rangle
    =
    c.
\]

After the cross-term coefficients are fixed, the verifier samples
\(y_0\ChooseRandom\Field^*\) and checks
Eq.~\eqref{eq:packed-ip-overview} at \(Y=y_0\).  This reduces the original
length-\(K\) claim to the full-field inner product
\[
    \left\langle
        \bm g_a(y_0),
        \bm g_b(y_0)
    \right\rangle
\]
of length \(K/\kappa\), which is then proved using EO-LiteBullet.

The prover computes \(\bm g_a(y_0)\) directly from the native \(b\)-bit
representation of \(\bm a\), while \(\bm g_b(y_0)\) is evaluated from the
compact tensor structure of \(\bm b\).  Consequently, the protocol
materializes neither a length-\(K\) full-field representation of
\(\bm a\) nor the expanded vector \(\bm b\).  Choosing
\(\kappa= \lfloor |\Field|/b \rfloor\) makes both the multiplication work and the
resident full-field state scale with the native input bit width.